\documentclass[11pt,a4paper]{article}
\usepackage[margin=2.3cm]{geometry}
\usepackage{amsmath,amssymb,booktabs,hyperref,enumitem}
\usepackage[T1]{fontenc}
\usepackage{xcolor}
\hypersetup{colorlinks=true,linkcolor=blue!60!black,urlcolor=blue!60!black}
\newcommand{\Q}{\operatorname{Q}}
\newcommand{\EbNo}{E_b/N_0}
\newcommand{\EbNJ}{E_b/N_J}
\setlength{\parskip}{4pt}

\title{Theory note S1 --- Frequency hopping under partial-band jamming:\\
processing gain, the worst-case band fraction, and the jamming margin}
\author{Chang Young Ham}
\date{September 2026 \quad (accompanies \texttt{src/s1\_fh\_partial\_band\_jamming.py})}

\begin{document}
\maketitle

\noindent\textbf{Purpose.}
S1 adds a hopping transmitter and an intelligent jammer to the S0 link. This
note derives (i) why hopping buys a processing gain, (ii) the closed-form BER
used to verify the simulator, (iii) the jammer's optimal band fraction
$\rho^\ast$ and the resulting inverse-linear BER law, and (iv) the jamming
margin numbers reported in \texttt{results/s1\_link\_budget.txt}.
Mathematics used: total probability, one-variable optimisation, a first-order
asymptotic of the Gaussian tail.

\section{Processing gain}

The data occupies bandwidth $W_d$; the transmitter hops among $N_{\mathrm{ch}}$
channels spread over $W_{ss}\approx N_{\mathrm{ch}}W_d$. A jammer with total
power $J$ that wants to cover the whole hopping band must spread $J$ over
$W_{ss}$, so its density is $N_J = J/W_{ss}$. The transmitter's signal at any
instant occupies only $W_d$, where the jammer delivers only
$N_J W_d = J\,W_d/W_{ss} = J/N_{\mathrm{ch}}$ of its power. The jammer
effectively loses a factor
\[
  G_p = \frac{W_{ss}}{W_d} \approx N_{\mathrm{ch}}, \qquad
  G_p\,[\mathrm{dB}] = 10\log_{10}N_{\mathrm{ch}}.
\]
With $N_{\mathrm{ch}}=100$: $G_p = 20$\,dB.

\paragraph{Relating $\EbNJ$ to $J/S$.}
With signal power $S$ and bit rate $R_b$, $E_b = S/R_b$. Taking $W_d\approx R_b$,
\[
  \frac{E_b}{N_J} = \frac{S/R_b}{J/W_{ss}} = \frac{S}{J}\cdot\frac{W_{ss}}{R_b}
  \approx \frac{S}{J}\,G_p
  \quad\Longrightarrow\quad
  \boxed{\ \frac{E_b}{N_J}\,[\mathrm{dB}] = G_p\,[\mathrm{dB}] - \frac{J}{S}\,[\mathrm{dB}]\ }.
\]
This is the line \texttt{ebnj\_lin = db\_to\_linear(gp\_db - js\_db)} in the code and
the reason the figure's $x$-axis is $J/S$ while the physics runs on $\EbNJ$.

\section{Partial-band jamming: the mixture formula}

An intelligent jammer does not spread over the whole band. It concentrates $J$
on a fraction $\rho\in(0,1]$ of the channels, raising its density there to
$N_J/\rho$ and leaving the rest clean. The transmitter does not know which
channels are jammed (pseudo-random pattern), so each hop lands in the jammed
region independently with probability $\rho$. By the law of total probability,
\begin{equation}
  P_b(\rho) = (1-\rho)\,\Q\!\left(\sqrt{\tfrac{2E_b}{N_0}}\right)
            + \rho\,\Q\!\left(\sqrt{\tfrac{2E_b}{N_0 + N_J/\rho}}\right).
  \label{eq:mix}
\end{equation}
The first term is the thermal-noise-only BER of S0 (clean hop); the second is an
S0 link whose noise density is $N_0$ plus the concentrated jammer density.
This is \texttt{ber\_partial\_band\_theory()} and it is what the Monte Carlo
(random jammed flag per hop, then the S0 machinery) is checked against.

\paragraph{Limit check ($\rho=1$).}
Eq.~\eqref{eq:mix} reduces to $\Q\!\big(\sqrt{2E_b/(N_0+N_J)}\big)$: plain AWGN with
the two densities added. \texttt{VERIFICATION 2} asserts exactly this
(agreement to $10^{-12}$).

\section{The jammer's optimal $\rho$ and the inverse-linear law}

From the jammer's point of view $\rho$ is a control knob and it wants to
maximise \eqref{eq:mix}. When the jammer dominates the thermal noise
($N_J/\rho \gg N_0$), drop $N_0$:
\[
  P_b(\rho) \approx \rho\,\Q\!\left(\sqrt{\tfrac{2\rho E_b}{N_J}}\right).
\]
Substitute $x = 2\rho\,\EbNJ$ (so $\rho = x/(2\EbNJ)$):
\[
  P_b \approx \frac{x\,\Q(\sqrt{x})}{2\,\EbNJ} = \frac{g(x)}{2\,\EbNJ},
  \qquad g(x) := x\,\Q(\sqrt x).
\]
The dependence on $\EbNJ$ has been pulled out; the jammer just maximises
$g(x)$ over $x>0$. Setting $g'(x)=0$, with $\Q'(u) = -\varphi(u)$ and
$\varphi(u)=e^{-u^2/2}/\sqrt{2\pi}$,
\[
  g'(x) = \Q(\sqrt x) - x\,\varphi(\sqrt x)\,\frac{1}{2\sqrt x}
        = \Q(\sqrt x) - \frac{\sqrt x}{2}\,\varphi(\sqrt x) = 0.
\]
This transcendental equation has one positive root, found numerically:
\[
  x^\ast \approx 1.418, \qquad g(x^\ast)\approx 0.1657 .
\]
\textbf{Stop here and check:} evaluate both sides of $\Q(\sqrt x)=\tfrac{\sqrt x}{2}\varphi(\sqrt x)$
at $x=1.418$; they agree to four decimals ($0.1169$).

Two consequences follow immediately.
\begin{align}
  \rho^\ast &= \frac{x^\ast}{2\,\EbNJ} \approx \frac{0.709}{\EbNJ}, \label{eq:rhostar}\\[2pt]
  P_b^{\text{worst}} &= \frac{g(x^\ast)}{2\,\EbNJ} \approx \frac{0.083}{\EbNJ}. \label{eq:worst}
\end{align}
Eq.~\eqref{eq:worst} is the central fact of S1: \emph{against an optimal
partial-band jammer the BER decays only as $1/(\EbNJ)$}, not exponentially as
in S0. Doubling the signal power halves the error rate instead of squaring it
away. (Eq.~\eqref{eq:rhostar} is only meaningful while $\rho^\ast\le 1$, i.e.\
$\EbNJ \ge x^\ast/2 \approx -1.5$\,dB; below that the jammer simply uses the
whole band.)

\paragraph{The number in the figure.}
At $J/S = 12$\,dB with $G_p=20$\,dB, $\EbNJ = 8$\,dB $=6.31$. Eq.~\eqref{eq:rhostar}
gives $\rho^\ast = 0.709/6.31 = 0.112$ with thermal noise ignored. Keeping the
thermal term $N_0$ in \eqref{eq:mix} shifts the optimum slightly upward; a
numerical search of the exact formula gives $\rho^\ast = 0.135$, which is the
$0.14$ reported by both the closed-form grid and the Monte Carlo
(\texttt{VERIFICATION 3}). The corresponding worst-case BER is
$1.42\times10^{-2}$ versus $2.70\times10^{-3}$ for a full-band jammer of the same
power: the factor of five in the figure.

\section{Jamming margin}

A link is declared \emph{up} if $P_b \le \mathrm{BER}_{\mathrm{out}} = 10^{-3}$
(assumed criterion; the S2 note discusses its consequences). Invert the BER
formula to find the $\EbNJ$ at which the criterion is just met,
$(\EbNJ)_{\mathrm{req}}$, then allow an implementation loss $L_{\mathrm{sys}}$.
The jamming margin is the $J/S$ the link can absorb:
\[
  M_J\,[\mathrm{dB}] = G_p - (\EbNJ)_{\mathrm{req}} - L_{\mathrm{sys}}.
\]
For $N_{\mathrm{ch}}=100$, $\EbNo=10$\,dB, $L_{\mathrm{sys}}=2$\,dB:

\begin{center}
\begin{tabular}{@{}lcc@{}}
\toprule
jammer strategy & $(\EbNJ)_{\mathrm{req}}$ & $M_J$ \\
\midrule
full band, $\rho=1$ & 9.65\,dB & $20-9.65-2 = \mathbf{8.35}$\,dB \\
optimal $\rho$      & 18.95\,dB & $20-18.95-2 = \mathbf{-0.95}$\,dB \\
\bottomrule
\end{tabular}
\end{center}

The optimal jammer needs 9.3\,dB \emph{less} power to break the link. The
\emph{negative} margin in the second row means that at $\mathrm{BER}_{\mathrm{out}}=10^{-3}$,
100 hop channels alone cannot protect a BPSK link against a jammer of equal
received power --- a statement that motivates every extension listed in S2.

\section{Correspondence with the code}

\begin{center}
\begin{tabular}{@{}ll@{}}
\toprule
Mathematics & Code \\
\midrule
$G_p = 10\log_{10}N_{\mathrm{ch}}$ & \texttt{gp\_db = linear\_to\_db(N\_CH)} \\
$\EbNJ = G_p - J/S$ & \texttt{ebnj\_lin = db\_to\_linear(gp\_db - js\_db)} \\
jammed hop w.p.\ $\rho$ & \texttt{jammed = rng.random(...) < rho} \\
$\sigma^2 = (N_0 + N_J/\rho)/2$ & \texttt{sigma\_jammed = np.sqrt((n0 + nj / rho) / 2.0)} \\
Eq.~\eqref{eq:mix} & \texttt{ber\_partial\_band\_theory()} \\
$\rho^\ast$ by grid search & \texttt{RHO\_GRID[np.argmax(ber\_rho\_theory)]} \\
$(\EbNJ)_{\mathrm{req}}$ & first grid point with worst-case BER $\le 10^{-3}$ \\
\bottomrule
\end{tabular}
\end{center}

\section{Self-check}
\begin{enumerate}[nosep]
  \item Explain in one sentence why a jammer that covers the whole hopping band
        wastes $1-1/N_{\mathrm{ch}}$ of its power.
  \item Derive \eqref{eq:mix} from the law of total probability.
  \item Starting from $P_b\approx\rho\,\Q(\sqrt{2\rho\EbNJ})$, show that the optimal
        $\rho$ scales as $1/(\EbNJ)$ and that the worst-case BER scales the same
        way. Why does this matter more than the exact constant $0.083$?
  \item Why is the required $\EbNJ$ for the optimal jammer (18.95\,dB) so much
        larger than for the full-band jammer (9.65\,dB)? Relate the gap to
        the $1/x$ versus $e^{-x}$ decay laws.
  \item What would change in the analysis if the jammer \emph{knew} the
        hopping pattern (a follower jammer)? Which assumption breaks?
\end{enumerate}

\end{document}
